%% ma: L11-B27-0021
%% lop: 11
%% chuong: 7
%% bai: 27
%% dang: 11.27.17
%% loai: TN4
%% muc: TH
%% dap_an: "A"
%% nguon: "(NBV)-ĐỀ SỐ 19-MỖI NGÀY 1 ĐỀ THI 2025.pdf (D19, MỖI NGÀY 1 ĐỀ - Đề số 19), Phần I Câu 9"
%% kien_thuc: "Bài 27 (thể tích khối chóp), hình chóp tứ giác đều"
%% trang_thai: chinh-thuc
%% ghi_chu: "Đáp án gốc A đúng (giải độc lập: h = 3V/S = 3·2592100/52900 ≈ 147,0 m). Bỏ ảnh minh họa kim tự tháp (ảnh trang trí, không có số liệu toán học). Đã duyệt 19/10/2026."
\begin{ex}
Kim tự tháp Kheops ở Ai Cập được xây dựng vào khoảng $2500$ năm trước Công nguyên. Kim tự tháp này là một khối chóp tứ giác đều có thể tích khoảng $2\,592\,100\ \mathrm{m^3}$, cạnh đáy dài $230$ m. Chiều cao của kim tự tháp là
\choice{\True $147$ m}{$145$ m}{$152$ m}{$150$ m}
\loigiai{Diện tích đáy là $S=230^2=52\,900\ \mathrm{m^2}$. Từ $V=\dfrac13Sh$ suy ra $h=\dfrac{3V}{S}=\dfrac{3\cdot2\,592\,100}{52\,900}\approx147\ \mathrm{(m)}$. Chọn A.}
\end{ex}
